\section{Benchmark Construction and Leakage Audit}
\label{app:benchmarks}

\subsection{Passages, concepts and future questions}

The primary author model is GPT-4.1. Each episode contains a passage,
candidate concepts with construction-assigned utility labels, and a
future question and answer. Explicit places useful content in the
passage. \textbf{Implied-answer} training episodes cue an unstated
useful concept that itself supplies the answer, including episodes from
the alternate-author benchmark. The extended tables use this name for
that training construction. Decoupled requires an unstated bridge to support a further
inference: the answer is absent from the passage and candidate set,
and the question shares no passage content word except a protagonist's
name. Compositional requires a multi-hop inference and permits question
anchors from the passage. Decoupled and Compositional are development
batteries used for configuration selection and excluded from gradient
training.

Mechanical validation checks candidate strings, bridge absence, answer
exclusion and the Decoupled question's lexical overlap. Concept
deduplication removes repeated concepts across retained episodes.
The leakage audit additionally checks full lexemes, inflected answers
and question anchors. A blind GPT-6-Astra alias audit receives the
future question, gold answer and bridge. It rejects identity,
morphological or derivational variants, synonyms, translations,
abbreviations and containment that trivially yields the answer.
For example, \emph{caterer/catering} is an alias, whereas
\emph{Madrid/Spanish} requires a further inference. The audit records
the verdict and exclusion reason for each episode.

\subsection{Episode-authoring prompt}
\label{app:authorprompt}
GPT-4.1 receives the system instruction, ``You are a careful dataset
designer for a cognitive-science experiment on LLM memory.'' The
following prose transcription preserves the generation constraints;
the output-format instruction requests one structured record per episode.
\begin{sdabox}{Episode-authoring prompt}
Design a batch of diverse test episodes about silent bridge reasoning
for a memory experiment. Each episode has a passage of four to eight
sentences about a named person; invent a distinctive first name,
different in every episode, and vary the domain widely. A question
asked later, out of context, refers to the episode only through the
person's name and requires an unstated bridge fact inferable from the
passage, followed by common general knowledge. Supply a short correct
answer of one to three words and five to seven candidate concepts,
each a common lowercase word with a utility category and a short reason.

Exactly one candidate is the useful bridge concept, and its word must
be absent from the passage. The question must share no passage content
word except the person's name: someone who never read the passage
must be unable to answer it. For a passage about Elena wandering the
Plaza Mayor, ask what language Elena should have brushed up on for her
trip, without naming the plaza in the question. The answer must differ
from every candidate and be absent from the passage; combining the
bridge with general knowledge yields the answer. For example, the
Plaza Mayor cues \emph{madrid}, which supports \emph{spanish}. A reader
remembering only the bridge plus the question must be able to answer.

Include two or three vivid, emotional or surprising distractors that
occur literally in the passage but are useless for the question, and
one or two neutral concepts that occur in the passage. Every candidate
must be a single common English lowercase word. Do not reuse concepts
or person names across episodes.
\end{sdabox}
Generation uses temperature 1.0 and batches of up to five episodes.
Mechanical filters and the alias audit apply after generation.

\subsection{Frozen batteries and replication benchmarks}
\label{app:fresh}

Fresh-1 and Fresh-2 use the Decoupled construction prompt and mechanical
filters, followed by deduplication and disjointness filtering. A context
match after lowercasing and trimming against a prior battery causes
exclusion. A bridge matching a construction-labeled training concept
also causes exclusion. Fresh-2 additionally excludes contexts appearing
in Fresh-1. Shared answer categories, such as languages and currencies,
are permitted. The retained batteries contain 316 and 363 episodes,
respectively. The alternate-author benchmark repeats the inferred-concept
construction with a different author model. The entity-bridge benchmark
places a named-entity bridge among scene-object distractors.

\subsection{Reconstructing the reported controls}
\label{app:pools}
\label{app:controls}

All pool interventions traverse the 363 held-out Fresh-2 episodes in
their frozen order and preserve the question, answer and original
utility labels. Added candidates receive negative utility labels.
Within an episode, original candidates precede additions, whose recorded
order is retained. The frozen control archive records the passages,
ordered additions, named candidates and context permutation; replay uses
these recorded draws. The standard pool contains 1,607 candidates,
4.43 per episode. Statedness uses the case-insensitive stem match
specified in \S\ref{sec:setup}.

\paragraph{A worked source episode.}
The first episode reads: ``Yara compared renewal quotes from different
insurers online and reviewed her accident history. She called the hotline
to ask clarifying questions about premium increases. After considering
the coverage limits and deductibles, Yara finally entered her payment
details to confirm the policy.'' Its candidates, in order, are
\emph{insurer, deductible, history, policy, insurance}; \emph{insurance}
is useful. The question is ``Which professional would be most relevant
for Yara to contact regarding her recent purchase?'' and the answer is
\emph{agent}. The second episode is Jibril's visit in
Figure~\ref{fig:episode}.

\paragraph{Stated-bridge, position and distractor-named controls.}
The exact sentence form is ``The passage also mentions \emph{concept}.''
The stated-bridge control appends this sentence with the useful concept;
the position control prepends it. The distractor-named control appends
it with the first non-useful candidate in the original order. Thus
Yara's stated-bridge control adds ``The passage also mentions insurance.''
after the passage, the position control puts that sentence before
``Yara compared,'' and the distractor-named control instead appends
``The passage also mentions insurer.'' These three controls keep all
five candidates and their order unchanged.

\paragraph{Pool with unstated distractors.}
An author model proposes four single-word concepts from the useful
concept's semantic category, asking for alternatives a careless reader
might confuse with it but that the passage does not imply. The instruction
requires lowercase words absent from the passage and a comma-separated
response. We retain the response order, discard a proposal that matches
an original candidate or occurs as a case-insensitive substring of the
passage, and append the accepted concepts. The frozen pool contains
1,450 additions: four in 361 episodes and three in episodes 24 and 263
after filtering, giving 3,057 candidates in total (8.42 per episode).
For Yara, the ordered additions are \emph{warranty, mortgage, pension,
annuity}; her passage and five original candidates stay fixed.
The development confusable-distractor control applies the same procedure
to Decoupled. The unrelated-distractor control instead takes bridges
from subsequent episodes, as described for the enlarged pool below.

\paragraph{Statedness-balanced pool.}
Odd positions, counting from one, keep the original passage; even
positions append ``The passage also mentions \emph{concept}.'' with
the useful concept substituted. Every episode keeps its original
candidates and takes the first four accepted additions at odd positions
or the first three at even positions from the pool with unstated
distractors, preserving their order. This yields 2,877 candidates,
7.93 per episode. Yara occupies position one, so her passage is unchanged
and her pool is the five original candidates followed by \emph{warranty,
mortgage, pension, annuity}. Jibril occupies position two: append
``The passage also mentions dubai.'' to Figure~\ref{fig:episode}'s
passage and append \emph{doha, riyadh, muscat} to its six candidates.
The original question and answer remain fixed in each case.

\paragraph{Naming-balanced pool.}
\label{app:namingpool}
Use the same ordered candidate lists as the statedness-balanced pool,
but start from the original passages. Name four candidates per episode.
At odd positions exclude the useful concept and sample four other
candidates uniformly without replacement; at even positions include the
useful concept and sample three others. Use one Python standard
pseudorandom generator seeded with zero, process episodes in frozen
order, prepend the useful concept to the sampled list when included,
and shuffle each selected list before constructing sentences. Sampling
and shuffling consume the same continuing random stream. The five
sentence forms, in rotation, are ``The passage also mentions
\emph{concept}.'', ``\emph{concept} is relevant here.'', ``Note that
\emph{concept} comes up as well.'', ``There is also \emph{concept} to
consider.'', and ``\emph{concept} features in this as well.''

Number episodes and insertions from zero for the placement rule. The
first selected candidate uses the sentence form at the episode's
position in the five-form cycle; each subsequent candidate advances
one form. Strip trailing full stops from the original passage, split
at full-stop--space boundaries, trim each segment and restore one full
stop to each. Before each insertion, count the current sentence list,
including earlier additions. The insertion boundary is the remainder
of the episode index plus three times the insertion index upon division
by that count plus one. Boundary zero precedes the passage; the final
boundary follows it. Insertions thus advance through an evolving
sentence list rather than all occurring at its end.

For Yara, the shuffled selection is \emph{pension, annuity, insurer,
policy}. The additions are ``The passage also mentions pension.'',
``annuity is relevant here.'', ``Note that insurer comes up as well.'',
and ``There is also policy to consider.'' Their successive insertion
boundaries are zero, three, zero and two. The final passage begins with
the insurer, pension and policy sentences, then Yara's first two original
sentences, then the annuity sentence, then her final original sentence.
There are 1,452 naming interventions, 181 of which name the useful
concept: a named candidate is useful 12.5\% of the time, against the
12.6\% base rate among 2,877 candidates. Questions, answers and
candidate order stay fixed.

\paragraph{Enlarged pool.}
Combine the standard candidate list, the accepted same-topic additions,
and four bridges from other episodes. To choose the latter, start with
the next episode in frozen order, advance cyclically, and accept a bridge
only if it is absent from the passage, absent from the original
candidate list and distinct from earlier accepted bridges. Stop after
four accepted bridges. Concatenate the three lists in this order and
retain the first occurrence of each concept. The resulting 4,506
candidates average 12.41 per episode. Yara's pool is \emph{insurer,
deductible, history, policy, insurance, warranty, mortgage, pension,
annuity, dubai, church, kanji, doctor}. Her passage, question and answer
remain unchanged. The unrelated-distractor control uses just the
standard candidates and the four cross-episode additions; for Yara
these are \emph{dubai, church, kanji, doctor}.

\paragraph{Shuffled context and bridge replacement.}
Shuffled context uses the archived episode permutation, replacing each
passage while preserving its own candidates, utility labels, question
and answer. The first episode receives the passage at position 105,
about Lars skating on slick tracks with thin-bladed boots; its candidates
remain Yara's five concepts and its answer remains \emph{agent}.
The bridge-replacement control acts on SDA's recorded top-two selections
only when they contain the useful concept: retain the other selection
and replace the useful concept with the next-ranked candidate, then
replay the same frozen reader and question. For a selection of
\emph{dubai, elevator} with \emph{gift} next, the replayed memory is
\emph{gift, elevator}; Jibril's question and \emph{arabic} answer stay
fixed. This example illustrates the replacement rule; the measured
249 replays use the recorded selections (\S\ref{app:episodemetrics}).

\paragraph{Candidate proposal and reader references.}
The context-only proposer receives the original passage and the
instruction in Appendix~\ref{app:prompts}; the future question and bridge
are hidden. Its parsed output replaces the supplied candidate list in
proposal order. For Yara the resulting candidates are \emph{online,
accident, hotline, premium, coverage, deductibles, payment, policy},
so \emph{insurance} is unavailable to admission. At two stored concepts,
the same frozen reader answers from each policy's selected pair.
The question-only reference gives that reader Yara's question with no
memory; the whole-passage reference gives the same question and her
original three-sentence passage. Neither reference performs admission.
The wording controls change only the admission instruction, whose five
exact forms appear in Appendix~\ref{app:wording}. Authorship and
entity-bridge controls repeat the episode construction with the authors
and concept categories specified in Appendix~\ref{app:replications}.

SDA improves ranking and retention under the naming-balanced construction
(Table~\ref{tab:namingbalanced}). Across Qwen3-8B seeds, within-episode
AUC is 0.715--0.725, versus 0.566 for the report and 0.636 for the
workspace readout. Useful-concept retention at $k=2$ is 53.2--54.0\%,
versus 24.0\% and 40.8\%, respectively.

\begin{table}[htbp]
\centering
\caption{Naming-balanced pool recall across 363 episodes. Every episode names four candidates using five sentence templates; a named candidate is useful at the base rate. Recall is reported in percent under blind semantic grading at one, two and three stored concepts for six models. Report is the model's own yes/no importance report; Rating is the 1--10 importance rating \citep{park2023generative}; SDA is the mean of three seeds.}
\label{tab:namingrecall}
\small
\setlength{\tabcolsep}{3.5pt}
\begin{tabular}{lccccccccc}
\toprule
& \multicolumn{3}{c}{$k=1$} & \multicolumn{3}{c}{$k=2$} & \multicolumn{3}{c}{$k=3$} \\
\cmidrule(lr){2-4} \cmidrule(lr){5-7} \cmidrule(lr){8-10}
Model & Report & Rating & \textbf{SDA} & Report & Rating & \textbf{SDA} & Report & Rating & \textbf{SDA} \\
\midrule
Qwen3-8B & 27.6 & 35.3 & \textbf{36.4} & 42.7 & 45.2 & \textbf{46.2} & 46.3 & 47.1 & \textbf{47.0} \\
Qwen3-14B & 21.8 & 30.3 & \textbf{33.7} & 33.3 & 36.4 & \textbf{40.8} & 41.3 & 43.2 & \textbf{41.2} \\
Qwen3-32B & 31.7 & 30.0 & \textbf{41.1} & 43.2 & 43.0 & \textbf{48.9} & 50.7 & 51.8 & \textbf{53.4} \\
Qwen2.5-7B & 9.4 & 11.6 & \textbf{13.3} & 21.2 & 25.6 & \textbf{28.3} & 28.6 & 33.3 & \textbf{33.0} \\
OLMo-2-7B & 24.8 & 30.6 & \textbf{37.0} & 39.4 & 43.0 & \textbf{46.7} & 44.6 & 52.1 & \textbf{49.5} \\
Mistral-7B & 4.7 & 5.8 & \textbf{9.9} & 9.6 & 13.2 & \textbf{16.0} & 19.8 & 21.5 & \textbf{23.7} \\
\bottomrule
\end{tabular}
\end{table}

\begin{table}[htbp]
\centering
\caption{Candidate ranking and useful-concept retention on the naming-balanced pool, Qwen3-8B, 363 episodes and 2,877 candidates. Every episode names four candidates using five sentence templates at varying positions. A named candidate is useful 12.5\% of the time, against a 12.6\% base rate. AUC uses construction utility labels; retention is in percent.}
\label{tab:namingbalanced}
\small
\setlength{\tabcolsep}{4pt}
\begin{tabular}{lccccc}
\toprule
Admission policy & \shortstack{Pooled\\AUC} & \shortstack{Episode\\AUC} & $B_1$ & $B_2$ & $B_3$ \\
\midrule
Yes/no importance report & 0.570 & 0.566 & 9.4 & 24.0 & 42.7 \\
Workspace readout & 0.635 & 0.636 & 22.0 & 40.8 & 54.8 \\
\textbf{SDA, seed 0} & \textbf{0.726} & \textbf{0.725} & \textbf{36.1} & \textbf{54.0} & \textbf{67.5} \\
\textbf{SDA, seed 1} & \textbf{0.714} & \textbf{0.715} & \textbf{32.2} & \textbf{54.0} & \textbf{65.6} \\
\textbf{SDA, seed 2} & \textbf{0.728} & \textbf{0.725} & \textbf{35.5} & \textbf{53.2} & \textbf{68.0} \\
\bottomrule
\end{tabular}
\end{table}


\section{SDA Training and Scoring Details}
\label{app:training}

\subsection{Teacher construction and data split}

Within each episode, let $r_S(c)$ be a candidate's ascending mid-rank
under readout $S$, with ties assigned the average of their occupied
ranks. For $m>1$ candidates, the percentile is
\begin{equation}
P_S(c;x)=\frac{r_S(c)-1}{m-1},
\qquad S\in\{W,\mathrm{LL}\}.
\label{eq:percentile}
\end{equation}
For a singleton pool the percentile is one half. Averaging the
percentiles yields the teacher in Eq.~\eqref{eq:teacher}; descending
teacher score and then candidate order determine the top-2 labels.
Both frozen readouts score the same episode and candidate before these
labels are constructed.

A source group is one author-and-construction battery: Explicit
(88 episodes), implied-answer from the primary author (75), or
implied-answer from the alternate author (57). The split is stratified
by these groups and assigns whole episodes, including all their candidates,
to training or validation. Within each group, episodes are sorted by a
cryptographic hash of the fixed split seed and the source-qualified
episode identity; the lowest 20\%, rounded up, enter validation.
The split seed is zero and stays fixed across optimization seeds.
This gives 70/18, 60/15 and 45/12 training/validation episodes in the
three groups, respectively: 175 training and 45 validation episodes in total. The teacher generates every
optimization target from these passages and candidates.
The selection protocol is specified in \S\ref{sec:student}.

\subsection{Complete training configuration}
\label{app:configuration}

LoRA adapts the attention query, key, value and output projections and
the feed-forward gate, up and down projections. Base parameters remain
frozen. Each model supplies its own teacher rankings and receives its
own admission adapter. The two highest-ranked candidates in each
training episode receive positive admission targets; all remaining
candidates receive negative targets. Only the target word contributes
to the vocabulary-normalized cross-entropy in Eq.~\eqref{eq:loss}.
Prompt tokens are masked. Qwen3 uses its chat template with thinking
disabled.

\begin{center}
\small
\begin{tabular}{@{}p{0.30\linewidth}p{0.64\linewidth}@{}}
\toprule
Setting & Configuration \\
\midrule
Adapter & Low-rank adaptation, rank 16, scaling 32; attention and
feed-forward projections \\
Optimizer & AdamW; constant learning rate $1\times10^{-5}$ \\
Batching & Batch size 1; gradient accumulation over 4 batches \\
Training duration & Two epochs; 414 optimizer steps \\
Sequence length & At most 1024 tokens \\
Data & 175 training and 45 validation episodes; source-stratified, disjoint episodes \\
Targets & Two highest teacher ranks positive; all others negative \\
\bottomrule
\end{tabular}
\end{center}

\subsection{Measurement conventions}
\label{app:conventions}

The yes/no report in every main table is scored in log space over the
yes and no token sets. The wording comparison uses its own common
scorer for all five wordings, so its rows are comparable to one another
under that aggregation: Mistral-7B's passage-only wording gives AUC
0.411 in the main tables and 0.440 in the wording comparison because
the two scorers aggregate differently; scoring the same wording one
prompt at a time also gives 0.440, identical to the batched result.
On Qwen3-8B's standard pool, first-token workspace AUC is 0.543 at
the end of the passage and 0.530 when the candidate is appended and
the state is read immediately before its first token; the latter is
the matched baseline for full-candidate AUC 0.574.
Frozen-report recall is 40.8\% in the main tables and 40.5\% in the
objective ablation, which evaluates the untrained checkpoint through
the adapter-loading path instead of the frozen path.

\subsection{Admission baselines}
\label{app:baselines}

SDA is compared with the following baselines. The \textbf{yes/no
importance report} uses the untrained model's score $V$ in
Eq.~\eqref{eq:score}, with Qwen3 thinking disabled.
The \textbf{1--10 importance rating} is the expected rating under
normalized probabilities of the complete answer strings.
The \textbf{reversed yes/no report} admits the lowest-scoring candidates.
\textbf{Admit what the passage omits} prioritizes candidates absent from
the passage under the statedness test and breaks ties at random.
\textbf{Omission then report} prioritizes unstated candidates and
uses the yes/no importance report to order candidates within each
provenance. The \textbf{provenance-adjusted report} adds a single
coefficient times the candidate's statedness indicator to the untrained
report, with the coefficient fitted on development episodes and held
fixed for evaluation. The \textbf{combined teacher used directly} ranks
by Eq.~\eqref{eq:teacher} without training an adapter.
The \textbf{implied-content wording} permits implied
content; the \textbf{provenance-cued wording} explicitly says the
passage may never mention the useful concept (App.~\ref{app:wording}).
\textbf{Random selection} samples candidates uniformly; \textbf{output
likelihood} and the \textbf{workspace readout} rank candidates directly
by their respective scores. Each baseline sees the passage and candidates
with the future question hidden. \textbf{Oracle labels} prioritize the
intended bridge; question-only and whole-passage readers provide recall
references. Bold entries identify SDA throughout the comparison tables.

\subsection{Exact admission, reader, rating and proposal instructions}
\label{app:prompts}

The boxes below reproduce the natural-language instructions. Italic
phrases stand for the supplied passage, candidate, memory list or
question. Each is a user message rendered with the model's chat template
and assistant generation prefix; Qwen3 thinking is disabled. SDA and the
report baseline share the admission prompt.

\begin{sdabox}{Admission prompt}
\emph{Passage}

Based only on the passage above, is the concept ``\emph{candidate
concept}'' one of the most important things to remember in order to
answer possible future questions? Answer with a single word: yes or no.
\end{sdabox}
The affirmative token set contains the distinct first tokens of
\emph{yes} in lowercase, initial-capital and uppercase forms, with
and without a leading space. The negative set is constructed analogously
from \emph{no}. Equation~\eqref{eq:score} gives the common report
score; SDA applies the same scoring rule with its adapter enabled.

\begin{sdabox}{Reader prompt}
Earlier you read a passage. The only things you remember from it are:
\emph{stored concepts separated by commas}.\\
Using only what you remember, answer concisely.\\
Question: \emph{future question}\\
Answer:
\end{sdabox}
The reader uses the frozen model with its admission adapter disabled.
The question-only reference asks ``Answer this question concisely.''
followed by the question and answer prefix. The whole-passage reference
supplies the passage, then ``Answer concisely.'' and the same question
and answer prefix. Greedy decoding is used throughout.

\begin{sdabox}{Rating baseline prompt}
\emph{Passage}

On a scale of 1 to 10, where 1 is purely mundane and 10 is extremely
important, rate how important it is to remember the concept
``\emph{candidate concept}'' in order to answer possible future questions
about this passage. Answer with a single number.
\end{sdabox}
Let $a_r$ be the complete answer string for rating $r$, and let $u_R(c,x)$
be this prompt. Each answer receives its full sequence probability,
multiplying the conditional probabilities of its tokens. The score is
\begin{equation}
\begin{split}
R(c;x)&=\sum_{r=1}^{10}r\,\pi_r(c;x),\\
\pi_r(c;x)&=\frac{p_{\mathcal M}(a_r\mid u_R(c,x))}
{\sum_{j=1}^{10}p_{\mathcal M}(a_j\mid u_R(c,x))}.
\end{split}
\label{eq:rating}
\end{equation}
A rating represented by multiple tokens is scored as a complete answer.

\begin{sdabox}{Candidate-proposal prompt}
\emph{Passage}

A memory system can store only a few single-word concepts about the
passage above. List up to 8 concepts that could be worth remembering
for possible future questions, including things the passage clearly
implies but does not state. Answer with lowercase single words separated
by commas, and nothing else.
\end{sdabox}
The context-only proposer uses Qwen3-8B with greedy decoding. Distinct
parsed concepts form the candidate pool. The bridge, future question
and answer are withheld from proposal and carried forward for evaluation.
Bridge coverage accepts case differences and a trailing plural ending.

\subsection{Ablation protocol}
\label{app:ablationprotocol}

The teacher and optimization ablations in \S\ref{sec:ablations} use the
same training episodes, adapter shape and matched step budget. The
recall-reward condition samples admitted sets from the admission scores
and asks the frozen reader to answer from the selected memory. Answer
correctness supplies the reward. Within each rollout group, reward
differences are normalized into advantages; updates use a clipped
group-relative policy objective and a penalty for departure from the
frozen policy \citep{shao2024deepseekmath}. The future question enters
the reward computation. The refinement condition applies these updates
after SDA training.

\section{Documented Prospective Freeze}
\label{app:freeze}

The dated experiment record documents prospective freezes of Fresh-1
and Fresh-2 on 15 September 2026. It records each battery's episode
count and cryptographic digest before evaluated models read it: 316
episodes for Fresh-1 and 363 for Fresh-2. The accompanying reproducibility
record retains the dated entries and verified battery digests.
Fresh-1 supplied the initial frozen-battery evaluation.
Fresh-2 was frozen before any evaluated model read it
and after configuration selection. Its primary endpoint compares SDA
with the yes/no report; a primary training seed was designated and the
remaining seeds reported alongside it. Its secondary endpoint compares
SDA with the rating baseline. Fresh confirmation uses the blind semantic
judge in App.~\ref{app:judge}.

\section{Recall Grading and Serving Protocol}
\label{app:repro}

\subsection{Blind semantic judge}
\label{app:judge}

The judge, GPT-6-Astra, receives the question, gold answer and reader
response with no admission-policy identity. The substantive grading
instruction is quoted verbatim; its output-format instruction requests
a label for every supplied answer.
\begin{sdabox}{Grading instruction}
Grade short answers to memory questions. For each item decide whether
ANSWER gives GOLD as its answer.\\
Labels:\\
``correct'': the answer commits to GOLD or an exact equivalent (synonym,
inflection, fuller name) as its single answer. Extra explanation is fine.\\
``hedged'': it offers two or more alternative answers (e.g.\ ``X or Y''),
or states the answer cannot be determined, even if GOLD is mentioned.\\
``incorrect'': anything else, including a different answer.
\end{sdabox}
Only correct responses contribute to recall; hedged and incorrect
responses count as errors. Judgments are reused for identical question,
gold answer and response triples. The harness additionally grades each
answer by string matching after normalization. Agreement is 85.1\%
across 13,068 answers; every training run exceeds its own report under
both graders, with maximum exact McNemar $p=0.026$.

\subsection{Amortized scoring time and throughput}
\label{app:timing}

The timing evaluation loads Qwen3-8B on one A100-80GB and constructs an
admission prompt for each of the 1,607 standard-pool Fresh-2 candidates. vLLM enables
prefix caching and LoRA serving, warms the engine on an initial subset,
and times batched generation over the full candidate list. Each prompt
generates one token greedily; SDA attaches the adapter to the admission
request, while the report baseline uses the same engine without it.
Model loading is outside the timed region. Milliseconds per candidate
are wall time divided by candidate count; throughput is candidate count
divided by wall time. The served and reference selections are passed
to the same frozen reader and semantic judge. The serving engine
recovers yes/no masses from returned token log probabilities.

\subsection{Capability with the adapter enabled on every token}
\label{app:capability}

Enabling the adapter on every token changes the four measured capability
scores by at most 0.3 percentage points under matched prompts.
Base and adapter scores are
71.6\% and 71.5\% on MMLU \citep{hendrycks2021mmlu}, 80.4\% and
80.7\% on ARC-Easy, 56.2\% and 56.4\% on ARC-Challenge
\citep{clark2018arc}, and 88.4\% and 88.6\% on the GSM8K evaluation
subset \citep{cobbe2021gsm8k}. The deployed interface attaches the
adapter to admission requests.


\section{Recall across Evaluation Batteries and Pools}
\label{app:recall}

SDA exceeds the report in all 54 comparisons across six
models, three seeds and three budgets. All seeds reach exact McNemar
$p<0.05$ in 17 of 18 model--budget cells; Qwen3-32B at $k=3$ has
maximum $p=0.073$. Across the six models, pooled report AUC is
0.658--0.736 and workspace AUC is 0.719--0.812.


\begin{table}[htbp]
\centering
\caption{Paired SDA minus baseline recall differences on the statedness-balanced pool, Qwen3-8B, for every row of Table~\ref{tab:balancedbaselines}, which cites each published method (1--10 importance rating \citep{park2023generative}; ConsistencyGate \citep{consistencygate2026}; A-MAC \citep{amac2026}; prediction error \citep{nemori2026}). Entries average three seeds and give 95\% intervals from 4,000 paired bootstrap draws over 363 episodes. Differences are in percentage points.}
\label{tab:balancedpaired}
\footnotesize
\setlength{\tabcolsep}{3pt}
\begin{tabular}{lccc}
\toprule
Baseline & $k=1$ & $k=2$ & $k=3$ \\
\midrule
\multicolumn{4}{l}{\emph{Importance judgments of the untrained model}} \\
Yes/no importance report & $+9.9$ [$+6.1$, $+13.8$] & $+10.7$ [$+7.0$, $+14.6$] & $+6.3$ [$+2.9$, $+9.7$] \\
1--10 importance rating & $+4.7$ [$+1.0$, $+8.3$] & $+5.2$ [$+1.7$, $+8.8$] & $+1.1$ [$-2.4$, $+4.6$] \\
Implied-content wording & $+9.1$ [$+5.4$, $+12.9$] & $+9.1$ [$+5.7$, $+12.7$] & $+4.1$ [$+1.0$, $+7.3$] \\
\addlinespace
\multicolumn{4}{l}{\emph{Admission criteria from published methods}} \\
ConsistencyGate, sampled & $+16.5$ [$+11.6$, $+21.5$] & $+18.7$ [$+14.1$, $+23.6$] & $+11.0$ [$+6.3$, $+15.5$] \\
ConsistencyGate, single pass & $+16.0$ [$+10.8$, $+21.1$] & $+18.7$ [$+14.0$, $+23.5$] & $+9.4$ [$+4.8$, $+14.1$] \\
A-MAC utility & $+10.5$ [$+5.6$, $+15.3$] & $+14.1$ [$+9.3$, $+19.0$] & $+6.9$ [$+2.6$, $+11.1$] \\
Prediction error & $+24.0$ [$+18.9$, $+29.1$] & $+27.8$ [$+22.5$, $+33.2$] & $+21.5$ [$+16.4$, $+26.6$] \\
\addlinespace
\multicolumn{4}{l}{\emph{Provenance rules}} \\
Provenance-adjusted report & $+3.9$ [$-0.9$, $+8.5$] & $+14.6$ [$+9.5$, $+19.8$] & $+11.8$ [$+6.4$, $+17.0$] \\
Reversed yes/no report & $+16.8$ [$+11.5$, $+22.2$] & $+24.0$ [$+18.5$, $+29.5$] & $+18.7$ [$+13.2$, $+24.3$] \\
Admit what the passage omits & $+11.8$ [$+6.4$, $+17.3$] & $+21.8$ [$+16.3$, $+27.2$] & $+15.4$ [$+10.0$, $+21.1$] \\
Omission then report & $+3.6$ [$-1.0$, $+8.6$] & $+14.6$ [$+9.4$, $+19.7$] & $+11.8$ [$+6.5$, $+17.2$] \\
\addlinespace
\multicolumn{4}{l}{\emph{Internal signals used directly}} \\
Output likelihood used directly & $+5.8$ [$+0.6$, $+10.7$] & $+11.8$ [$+6.7$, $+16.9$] & $+10.2$ [$+5.3$, $+15.1$] \\
Combined teacher used directly & $+3.6$ [$-1.2$, $+8.2$] & $+8.8$ [$+3.9$, $+13.9$] & $+5.5$ [$+0.9$, $+10.0$] \\
\bottomrule
\end{tabular}
\end{table}


Qwen3-8B SDA improves mean recall over the report on each development
and frozen battery in Table~\ref{tab:sdamain}. The gains also occur
with unstated distractors and in the enlarged pool
(Tables~\ref{tab:stress} and~\ref{tab:bigpool}). Across all six models,
mean SDA recall exceeds the report and rating at every standard-pool
budget (Tables~\ref{tab:sdabudget} and~\ref{tab:baselines}).

\begin{table*}[t]
\centering
\caption{Recall accuracy (\%) at $k{=}2$ on Qwen3-8B, blind semantic grading.
Development batteries (left); two frozen fresh Decoupled batteries
(right). SDA: mean\,$\pm$\,SD over three seeds. $^{*}$every seed
exceeds the yes/no report with exact McNemar $p<0.05$. Rating: 1--10
importance rating \citep{park2023generative}.}
\label{tab:sdamain}
\small
\setlength{\tabcolsep}{3pt}
\begin{tabular}{lcccccc}
\toprule
 & \multicolumn{4}{c}{Development} & \multicolumn{2}{c}{Fresh (held out)} \\
\cmidrule(lr){2-5} \cmidrule(lr){6-7}
Admission policy & Decoupled & Compositional & Confusable & Unrelated & Fresh-1 & Fresh-2 \\
 & $n{=}68$ & $n{=}52$ & $n{=}68$ & $n{=}68$ & $n{=}316$ & $n{=}363$ \\
\midrule
Yes/no importance report & 35.3 & 46.2 & 35.3 & 36.8 & 34.8 & 40.8 \\
1--10 importance rating & 47.1 & 53.8 & 48.5 & 48.5 & 43.7 & 45.7 \\
\wsrow \textbf{SDA} & \textbf{49.5}\,$\pm$\,2.2 & \textbf{54.5}\,$\pm$\,1.1 & \textbf{48.5}\,$\pm$\,2.5 & \textbf{50.0}\,$\pm$\,2.5 & \textbf{44.2}\,$\pm$\,0.7$^{*}$ & \textbf{50.0}\,$\pm$\,1.4$^{*}$ \\
\addlinespace
\refrow Oracle labels & 51.5 & 65.4 & 51.5 & 51.5 & 50.0 & 54.3 \\
\bottomrule
\end{tabular}
\end{table*}


\begin{table}[tbp]
\centering
\caption{Recall on the pool with unstated distractors (Fresh-2 with four extra candidates per episode that the passage also never states, 8.4 candidates per episode). Blind semantic grading, recall in percent, three training seeds per model; $p$ bounds the exact McNemar test against each model's own yes/no report across those seeds.}
\label{tab:stress}
\small
\setlength{\tabcolsep}{4pt}
\begin{tabular}{lcccc}
\toprule
Model & Yes/no report & \textbf{SDA (3 seeds)} & $p$ vs.\ report & Oracle labels \\
\midrule
Qwen3-8B & 41.0 & \textbf{51.0 / 51.2 / 49.9} & $\leq$ $10^{-3}$ & 54.3 \\
Qwen3-14B & 32.2 & \textbf{40.8 / 39.4 / 39.9} & $\leq$ $10^{-3}$ & 46.6 \\
Qwen3-32B & 41.0 & \textbf{51.0 / 50.7 / 50.1} & $\leq$ $10^{-4}$ & 54.8 \\
Qwen2.5-7B & 23.7 & \textbf{31.4 / 28.4 / 30.6} & $\leq$ .019 & 31.7 \\
OLMo-2-7B & 38.8 & \textbf{52.3 / 51.0 / 51.2} & $\leq$ $10^{-5}$ & 54.8 \\
Mistral-7B & 12.9 & \textbf{19.0 / 19.6 / 17.9} & $\leq$ .015 & 16.8 \\
\bottomrule
\end{tabular}

\end{table}

\begin{table}[tbp]
\centering
\caption{Recall on the enlarged pool: Fresh-2 with 12.41 candidates per episode, up from 4.43. Qwen3-8B, blind semantic grading, recall in percent. Output likelihood is used directly. SDA entries give three training seeds; $p$ bounds the exact McNemar test against the report across seeds. Oracle labels prioritize the intended bridge.}
\label{tab:bigpool}
\small
\setlength{\tabcolsep}{3pt}
\begin{tabular}{lcccccc}
\toprule
Budget & \shortstack{Yes/no\\report} & \shortstack{Output\\likelihood} & \shortstack{Workspace\\readout} & \textbf{SDA (3 seeds)} & $p$ vs.\ report & \shortstack{Oracle\\labels} \\
\midrule
$k=1$ & 29.5 & 28.1 & 24.0 & \textbf{38.0 / 38.0 / 39.7} & $\leq$ $10^{-4}$ & 47.1 \\
$k=2$ & 40.5 & 32.0 & 34.7 & \textbf{50.7 / 52.6 / 51.2} & $\leq$ $10^{-5}$ & 54.3 \\
$k=3$ & 41.9 & 33.9 & 38.6 & \textbf{52.9 / 52.1 / 52.1} & $\leq$ $10^{-5}$ & 51.5 \\
\bottomrule
\end{tabular}

\end{table}


\begin{table}[t]
\centering
\caption{Memory budget across models. Recall accuracy (\%) on the held-out Fresh-2 battery (363 episodes) with $k$ stored concepts; each model admits memories with its own adapter and answers with its own frozen weights; blind semantic grading; Rating is the 1--10 importance rating \citep{park2023generative}; SDA is the mean over three seeds; $^{*}$every seed exceeds the yes/no report of the same model at that budget (exact McNemar $p<0.05$).}
\label{tab:sdabudget}
\small
\setlength{\tabcolsep}{4pt}
\begin{tabular}{lccccccccc}
\toprule
 & \multicolumn{3}{c}{$k{=}1$} & \multicolumn{3}{c}{$k{=}2$} & \multicolumn{3}{c}{$k{=}3$} \\
\cmidrule(lr){2-4} \cmidrule(lr){5-7} \cmidrule(lr){8-10}
Model & Report & Rating & \textbf{SDA} & Report & Rating & \textbf{SDA} & Report & Rating & \textbf{SDA} \\
\midrule
Qwen3-8B & 28.1 & 34.7 & \textbf{37.6}$^{*}$ & 40.8 & 45.7 & \textbf{50.0}$^{*}$ & 41.6 & 47.9 & \textbf{52.2}$^{*}$ \\
Qwen3-14B & 23.4 & 31.1 & \textbf{33.5}$^{*}$ & 32.5 & 39.1 & \textbf{40.5}$^{*}$ & 37.5 & 43.5 & \textbf{43.6}$^{*}$ \\
Qwen3-32B & 32.5 & 32.0 & \textbf{41.4}$^{*}$ & 41.0 & 44.6 & \textbf{50.2}$^{*}$ & 46.0 & 50.7 & \textbf{53.0}$^{*}$ \\
Qwen2.5-7B-Instruct & 9.6 & 9.6 & \textbf{16.8}$^{*}$ & 23.1 & 22.6 & \textbf{29.8}$^{*}$ & 29.5 & 35.0 & \textbf{35.8}$^{*}$ \\
OLMo-2-7B-Instruct & 27.5 & 33.1 & \textbf{48.5}$^{*}$ & 38.0 & 44.1 & \textbf{55.4}$^{*}$ & 44.4 & 52.6 & \textbf{56.9}$^{*}$ \\
Mistral-7B-Instruct & 5.2 & 9.6 & \textbf{12.4}$^{*}$ & 12.7 & 15.7 & \textbf{19.8}$^{*}$ & 19.3 & 22.0 & \textbf{24.2}$^{*}$ \\
\bottomrule
\end{tabular}
\end{table}


\begin{table*}[t]
\centering
\caption{Admission policies on Qwen3-8B: standard pool and pool with unstated
distractors. Blind semantic grading, recall in percent. Output
likelihood is used directly. Rows above the rule are untrained
baselines; dashes denote unmeasured settings.}
\label{tab:baselines}
\small
\setlength{\tabcolsep}{4pt}
\begin{tabular}{lcccccc}
\toprule
 & \multicolumn{3}{c}{Standard pool} & \multicolumn{3}{c}{\shortstack{Pool with unstated\\distractors}} \\
\cmidrule(lr){2-4}\cmidrule(lr){5-7}
Admission policy & $k=1$ & $k=2$ & $k=3$ & $k=1$ & $k=2$ & $k=3$ \\
\midrule
Yes/no importance report & 28.1 & 40.8 & 41.6 & 28.6 & 41.0 & 41.3 \\
1--10 importance rating \citep{park2023generative} & 34.7 & 45.7 & 47.9 & 35.8 & 46.6 & 48.8 \\
Reversed yes/no report & 35.5 & 46.6 & 49.0 & 23.4 & 28.9 & 30.9 \\
Admit what the passage omits & 46.6 & 53.4 & 49.6 & 31.7 & 35.3 & 36.9 \\
Random selection & 28.6 & 44.6 & 46.6 & --- & --- & --- \\
Output likelihood & --- & 48.8 & --- & 30.6 & 39.1 & 40.5 \\
Workspace readout & 29.8 & 44.6 & 49.9 & 26.7 & 39.4 & 44.9 \\
\midrule
\textbf{SDA (mean of 3 seeds)} & \textbf{37.6} & \textbf{50.0} & \textbf{52.2} & \textbf{37.6} & \textbf{50.7} & \textbf{51.6} \\
\addlinespace
\refrow Oracle labels & 47.1 & 54.3 & 51.5 & 47.1 & 54.3 & 51.5 \\
\bottomrule
\end{tabular}

\end{table*}


\begin{table*}[t]
\centering
\caption{Paired SDA--baseline recall differences on three candidate pools,
Qwen3-8B. Entries are seed-averaged differences in points with 95\%
intervals from 4,000 paired bootstrap draws over episodes.
All comparisons use the same frozen reader and blind semantic judge.}
\label{tab:paired}
\small
\setlength{\tabcolsep}{3pt}
\begin{tabular}{lccc}
\toprule
Baseline & $k=1$ & $k=2$ & $k=3$ \\
\midrule
\multicolumn{4}{l}{\emph{Standard pool}} \\
\quad Yes/no importance report & $+9.6$ [$+5.8$, $+13.3$] & $+9.3$ [$+5.2$, $+13.2$] & $+10.6$ [$+6.8$, $+14.1$] \\
\quad 1--10 importance rating & $+2.9$ [$-0.6$, $+6.4$] & $+4.3$ [$+0.9$, $+7.8$] & $+4.2$ [$+1.3$, $+7.2$] \\
\quad Provenance-cued wording & $+5.4$ [$+1.4$, $+9.4$] & $+7.1$ [$+2.9$, $+11.4$] & $+6.2$ [$+2.8$, $+9.7$] \\
\quad Reversed yes/no report & $+2.1$ [$-3.3$, $+7.4$] & $+3.5$ [$-1.9$, $+8.9$] & $+3.1$ [$-1.3$, $+7.5$] \\
\quad Admit what the passage omits & $-8.9$ [$-13.1$, $-4.6$] & $-3.4$ [$-7.7$, $+1.0$] & $+2.6$ [$-1.2$, $+6.2$] \\
\quad Omission then report & $-9.5$ [$-13.7$, $-5.3$] & $-3.4$ [$-7.0$, $+0.4$] & $-0.5$ [$-3.2$, $+2.4$] \\
\quad Random selection & $+9.0$ [$+4.1$, $+14.1$] & $+5.4$ [$+0.6$, $+10.3$] & $+5.6$ [$+1.4$, $+9.8$] \\
\addlinespace
\multicolumn{4}{l}{\emph{Pool with unstated distractors}} \\
\quad Yes/no importance report & $+8.9$ [$+5.1$, $+12.8$] & $+9.6$ [$+5.5$, $+13.8$] & $+10.3$ [$+6.3$, $+14.1$] \\
\quad 1--10 importance rating & $+1.7$ [$-1.6$, $+5.0$] & $+4.1$ [$+0.6$, $+7.6$] & $+2.9$ [$-0.4$, $+6.2$] \\
\quad Reversed yes/no report & $+14.1$ [$+9.0$, $+19.5$] & $+21.8$ [$+16.2$, $+27.6$] & $+20.8$ [$+15.3$, $+26.1$] \\
\quad Admit what the passage omits & $+5.9$ [$+0.6$, $+10.9$] & $+15.4$ [$+10.3$, $+20.8$] & $+14.7$ [$+9.6$, $+19.8$] \\
\quad Omission then report & $-8.4$ [$-12.7$, $-4.2$] & $+5.8$ [$+0.8$, $+10.8$] & $+8.9$ [$+4.0$, $+13.8$] \\
\quad Output likelihood & $+7.0$ [$+1.7$, $+12.0$] & $+11.6$ [$+6.5$, $+16.8$] & $+11.1$ [$+6.2$, $+16.1$] \\
\addlinespace
\multicolumn{4}{l}{\emph{Enlarged pool}} \\
\quad Yes/no importance report & $+9.1$ [$+5.2$, $+12.9$] & $+11.0$ [$+6.7$, $+15.2$] & $+10.5$ [$+6.6$, $+14.2$] \\
\quad 1--10 importance rating & $+2.2$ [$-1.3$, $+5.6$] & $+4.7$ [$+0.9$, $+8.4$] & $+4.1$ [$+0.9$, $+7.6$] \\
\quad Omission then report & $-5.5$ [$-9.6$, $-1.5$] & $+6.6$ [$+1.6$, $+11.6$] & $+12.1$ [$+7.1$, $+17.2$] \\
\quad Output likelihood & $+10.5$ [$+5.0$, $+15.7$] & $+19.6$ [$+14.3$, $+24.9$] & $+18.5$ [$+13.3$, $+23.6$] \\
\quad Random selection & $+17.1$ [$+12.4$, $+22.0$] & $+24.5$ [$+19.1$, $+29.9$] & $+22.6$ [$+17.4$, $+27.7$] \\
\bottomrule
\end{tabular}

\end{table*}


\subsection{Candidate-pool comparisons}
\label{app:poolresults}

\paragraph{The standard pool rewards omission.}
The useful concept is the only candidate the passage never states in
95\% of standard-pool episodes. The omission baseline therefore reaches 46.6\%,
53.4\% and 49.6\% recall at $k=1,2,3$, respectively; omission then
report reaches 47.1\%, 53.4\% and 52.6\%, and the reversed
report reaches 35.5\%, 46.6\% and 49.0\%. SDA's corresponding means
are 37.6\%, 50.0\% and 52.2\%. The omission cue identifies the useful
concept directly in these episodes.

\paragraph{Unstated distractors separate useful admission from omission.}
The pool with unstated distractors adds four unstated distractors per episode.
Under blind semantic grading, SDA exceeds each model's own yes/no
report in 54 of 54 comparisons across six models, three seeds and
budgets $k=1,2,3$. All three
seeds reach exact McNemar $p<0.05$ in 17 of 18 model--budget cells;
the exception is Mistral-7B at $k=3$. At $k=2$, the largest $p$ across
models and seeds is 0.019 (Table~\ref{tab:stress}).
On Qwen3-8B, direct output likelihood reaches 30.6/39.1/40.5\%
recall at $k=1,2,3$, against SDA's 37.6/50.7/51.6\% means.
SDA leads the measured non-oracle baselines at $k=2,3$;
omission then report leads at $k=1$
(Tables~\ref{tab:baselines} and~\ref{tab:paired}). At $k=2$, the omission baseline falls from
53.4\% to 35.3\% and the reversed report from 46.6\% to 28.9\%, both
below the report's 41.0\%; SDA reaches 51.0\%, 51.2\% and
49.9\%. Every SDA seed exceeds the omission baseline
at every budget, with exact McNemar $p\leq0.097$ at $k=1$ and
$p\leq6.5\times10^{-7}$ at $k=2,3$; every seed exceeds the reversed
report with $p\leq4.3\times10^{-7}$ across budgets. SDA recalls are
37.7\% / 36.4\% / 38.6\% at $k=1$ and 51.8\% / 51.5\% / 51.5\%
at $k=3$; the rating baseline reaches 35.8\% and 48.8\%.

\paragraph{Naming the bridge and changing its category test reversal.}
On the stated-bridge control at $k=2$, SDA reaches 55.4\%, 55.6\% and
52.9\% recall, against 44.9\% for the report and 39.9\% for the
reversed report. On the entity-bridge benchmark, the reversed report
reaches 0.281 admission AUC, while SDA reaches 0.987, 0.980 and 0.987,
with within-episode AUC 1.000. The direction required by reversal is a
property of candidate construction, and is opposite on this benchmark
(\S\ref{sec:gap}).

\paragraph{The enlarged pool preserves gains at two and three concepts.}
With 12.41 candidates per episode, SDA reaches 38.0\% / 38.0\% / 39.7\%
at $k=1$, 50.7\% / 52.6\% / 51.2\% at $k=2$, and
52.9\% / 52.1\% / 52.1\% at $k=3$ (Table~\ref{tab:bigpool}).
The rating baseline reaches 36.4\%, 46.8\% and 48.2\%, respectively;
every SDA seed leads it at every budget. The paired mean differences
are $+2.2$ points ($[-1.3,+5.6]$), $+4.7$ ($[+0.9,+8.4]$), and
$+4.1$ ($[+0.9,+7.6]$), respectively (Table~\ref{tab:paired}).
Omission then report has the highest recall at $k=1$; SDA leads the
measured non-oracle baselines at $k=2,3$. Direct output likelihood reaches
28.1/32.0/33.9\%, against SDA's 38.6/51.5/52.4\% means.
Random selection reaches 21.5\%, 27.0\% and
29.8\%, the report 29.5\%, 40.5\% and 41.9\%, and the workspace
readout 24.0\%, 34.7\% and 38.6\%.


\subsection{Within-episode ranking and bridge retention}
\label{app:episodemetrics}

SDA improves both within-episode ranking and retention of the intended
bridge at $k=2$ (Table~\ref{tab:episodemetrics}). Paired whole-episode
bootstrap intervals compare each score with the report baseline.
Conditional recall distinguishes retaining the bridge from answering
the question: for the primary SDA seed, recall is 54.6\% with the
bridge and 41.2\% without it, versus 48.2\% and 37.7\% for the report.

The paired replacement experiment replays the primary SDA seed's
two selections. On the 249 episodes whose memory contains the bridge,
we replace only that concept with SDA's next-ranked candidate, keeping
the other selection, question and frozen reader fixed. Recall changes
from 54.6\% to 37.8\%, a 16.9-point decrease: 59 previously correct
answers become incorrect and 17 become correct (exact McNemar
$p=1.4\times10^{-6}$). Retaining the bridge supplies information
used by the reader.

\begin{table}[tbp]
\centering
\caption{Within-episode admission ranking on the standard Fresh-2 pool (Qwen3-8B). AUC and its difference from the report have paired 95\% whole-episode bootstrap intervals (4,000 draws). Retention is the percentage of episodes whose two admitted concepts contain the useful concept.}
\label{tab:episodemetrics}
\small
\setlength{\tabcolsep}{4pt}
\begin{tabular}{lccc}
\toprule
Admission score & \shortstack{Per-episode AUC\\{}[95\% CI]} & \shortstack{$\Delta$ vs.\ report\\{}[95\% CI]} & \shortstack{Bridge retained\\($k=2$)} \\
\midrule
Yes/no report & 0.326 [0.286, 0.366] & \emph{reference} & 29.8 \\
1--10 rating \citep{park2023generative} & 0.608 [0.567, 0.648] & $+0.282$ [$+0.246$, $+0.320$] & 61.2 \\
Workspace readout & 0.547 [0.511, 0.584] & $+0.221$ [$+0.172$, $+0.271$] & 51.0 \\
Output likelihood & 0.619 [0.582, 0.654] & $+0.293$ [$+0.242$, $+0.344$] & 59.2 \\
\textbf{SDA (seed 0)} & \textbf{0.682} [0.641, 0.721] & $+0.356$ [$+0.320$, $+0.393$] & \textbf{68.6} \\
\textbf{SDA (seed 1)} & \textbf{0.679} [0.639, 0.718] & $+0.353$ [$+0.316$, $+0.391$] & \textbf{69.7} \\
\textbf{SDA (seed 2)} & \textbf{0.688} [0.648, 0.727] & $+0.362$ [$+0.325$, $+0.401$] & \textbf{69.2} \\
\bottomrule
\end{tabular}

\end{table}


\paragraph{Bridge labels and reader success measure different outcomes.}
Qwen3-8B recall is 22.9\% with the question alone, 63.9\% with the
whole passage, 54.3\% with oracle labels and 50.0\% with SDA.
Oracle labels prioritize the intended bridge; the reader answers from
the complete stored set and its own knowledge. On OLMo-2-7B and
Mistral-7B, SDA reaches 55.4\% and 19.8\%, versus 54.8\% and 16.8\%
for oracle labels.



\section{Provenance, Wording and Pool-Composition Measurements}
\label{app:provenance}

\subsection{Ranking among candidates with the same provenance}
\label{app:withinprovenance}

The report discriminates usefulness among unstated candidates.
Restricting the pool with unstated distractors to episodes with an
unstated useful concept and at least one unstated negative leaves
345 episodes, with 5.01 qualifying candidates per episode on average.
Within this group, report AUC is 0.880--0.897 across six models and
workspace AUC is 0.592--0.697. Every paired 95\% interval for workspace
minus report lies below zero (Table~\ref{tab:withinunstated}).
This measurement holds provenance fixed; the full-pool comparison
measures the competition between stated and unstated candidates.

\begin{table*}[t]
\centering
\caption{Useful-concept ranking restricted to unstated candidates in the pool with unstated distractors. The 345 qualifying episodes have an unstated useful concept and at least one unstated negative, averaging 5.01 candidates in the restricted set. Entries are pooled AUC; $W-V$ gives the workspace--report difference with a paired 95\% whole-episode bootstrap interval from 4,000 draws.}
\label{tab:withinunstated}
\small
\setlength{\tabcolsep}{5pt}
\begin{tabular}{lccc}
\toprule
Model & Yes/no report & Workspace readout & $W-V$ [95\% CI] \\
\midrule
Qwen3-8B & 0.881 & 0.639 & $-0.242$ [$-0.274$, $-0.211$] \\
Qwen3-14B & 0.880 & 0.654 & $-0.226$ [$-0.257$, $-0.196$] \\
Qwen3-32B & 0.896 & 0.697 & $-0.199$ [$-0.228$, $-0.170$] \\
Qwen2.5-7B & 0.897 & 0.592 & $-0.305$ [$-0.340$, $-0.270$] \\
OLMo-2-7B & 0.883 & 0.661 & $-0.221$ [$-0.252$, $-0.191$] \\
Mistral-7B & 0.880 & 0.672 & $-0.208$ [$-0.238$, $-0.178$] \\
\bottomrule
\end{tabular}
\end{table*}


\subsection{The full-pool AUC changes with the comparison pairs}
\label{app:composition}

On Fresh-2, the useful concept is stated in 5\% of episodes and
negative candidates in 99\%; in 95\% of episodes, the useful concept
is the sole candidate absent from the passage. Adding unstated
distractors leaves every original Qwen3-8B report score unchanged.
It adds 1,450 negative candidates to the original 1,244, while retaining
the same 363 positives. Useful-concept AUC against the added negatives
is 0.889, compared with 0.322 against the original negatives.
The pooled AUC is therefore the weighted average
\begin{equation}
\begin{split}
\mathrm{AUC}_{\mathrm{expanded}}
&=\frac{1244}{2694}\,\mathrm{AUC}_{\mathrm{original}}\\
&\quad+\frac{1450}{2694}\,\mathrm{AUC}_{\mathrm{added}}\\
&=0.627.
\end{split}
\label{eq:pooldecomposition}
\end{equation}
The added comparisons account for 53.8\% of all positive--negative
pairs. They raise AUC without changing the ordering of the original
candidates. At a two-concept budget, the report continues to prioritize
passage content: recall is 41.0\%, against SDA's 51.0\%, 51.2\% and
49.9\% across seeds (Table~\ref{tab:baselines}).

\subsection{Wording changes the importance ordering}
\label{app:wording}

Five matched wordings keep the episode, candidates and utility labels
fixed. They ask about passage-only importance, remove that restriction,
allow implied content (\emph{implied-content wording}), state that the
passage may omit the useful concept (\emph{provenance-cued wording}),
or specify a two-concept budget. The measurements cover four models
on the standard pool and the pool with unstated distractors
(Table~\ref{tab:wording}).
On Qwen3-32B, the provenance-cued wording raises standard-pool AUC
from 0.398 to 0.876; the corresponding Qwen3-8B values are 0.321 and
0.537. On Mistral-7B, the implied-content wording gives
0.518, compared with 0.440 for passage-only importance. The direction
and magnitude of the wording effect vary across models.

The provenance-cued wording's ranking changes when the useful concept
is stated. On Qwen3-8B, it reaches AUC 0.417 when the
useful concept is named and 0.621 on Explicit. Its standard-pool recall
is 32.2\%, 43.0\% and 46.0\% at $k=1,2,3$
(Table~\ref{tab:paired}).

The strongest Qwen3-8B prompt baseline among wordings that permit
inferences without specifying the useful concept's provenance is the
implied-content wording. Its recall rises from 28.1/40.8/41.6\% to
29.2/41.3/44.9\%. SDA trained and evaluated with this same wording
reaches 37.1/50.7/51.0\% (Table~\ref{tab:wordingdistill}); all seeds
exceed the wording alone at every budget, with exact McNemar
$p\leq6.9\times10^{-5}$. The passage-only wording SDA means are
37.6/50.0/52.2\% (Table~\ref{tab:baselines}). Training through the
implied-content wording adds recall beyond the wording change alone.

\begin{table*}[t]
\centering
\caption{Implied-content wording and distillation on the standard pool,
Qwen3-8B, 363 episodes. Recall is in percent. SDA averages three
seeds; the final row gives seed-averaged paired differences and
95\% bootstrap intervals from 4,000 draws.}
\label{tab:wordingdistill}
\small
\setlength{\tabcolsep}{4pt}
\begin{tabular}{lccc}
\toprule
Admission policy & $k=1$ & $k=2$ & $k=3$ \\
\midrule
Yes/no importance report & 28.1 & 40.8 & 41.6 \\
Implied-content wording & 29.2 & 41.3 & 44.9 \\
\textbf{SDA with implied-content wording} & \textbf{37.1} & \textbf{50.7} & \textbf{51.0} \\
SDA minus implied-content wording & $+7.9$ [$+4.5$, $+11.4$] & $+9.4$ [$+5.8$, $+13.1$] & $+6.1$ [$+3.5$, $+8.9$] \\
\bottomrule
\end{tabular}

\end{table*}


\begin{table*}[t]
\centering
\caption{Importance wording controls on 363 Fresh-2 episodes per pool. Entries are pooled useful-concept AUC, using the untrained model's yes/no log odds. All wordings and models are shown.}
\label{tab:wording}
\small
\setlength{\tabcolsep}{4pt}
\begin{tabular}{llcc}
\toprule
Model & Importance wording & Standard pool & \shortstack{Pool with unstated\\distractors} \\
\midrule
Qwen3-8B & Passage only & 0.321 & 0.628 \\
 & Unrestricted & 0.391 & 0.663 \\
 & Implied-content wording & 0.462 & 0.692 \\
 & Provenance-cued wording & 0.537 & 0.726 \\
 & Two-concept budget & 0.235 & 0.591 \\
\addlinespace
Qwen3-32B & Passage only & 0.398 & 0.671 \\
 & Unrestricted & 0.555 & 0.749 \\
 & Implied-content wording & 0.656 & 0.802 \\
 & Provenance-cued wording & 0.876 & 0.892 \\
 & Two-concept budget & 0.327 & 0.633 \\
\addlinespace
OLMo-2-7B & Passage only & 0.327 & 0.630 \\
 & Unrestricted & 0.358 & 0.643 \\
 & Implied-content wording & 0.420 & 0.675 \\
 & Provenance-cued wording & 0.547 & 0.721 \\
 & Two-concept budget & 0.301 & 0.614 \\
\addlinespace
Mistral-7B & Passage only & 0.440 & 0.680 \\
 & Unrestricted & 0.462 & 0.684 \\
 & Implied-content wording & 0.518 & 0.716 \\
 & Provenance-cued wording & 0.394 & 0.639 \\
 & Two-concept budget & 0.336 & 0.626 \\
\bottomrule
\end{tabular}

\end{table*}


For wording $j$ with prompt $u_j(c,x)$, the scorer is
\begin{equation}
\begin{split}
s_j(c;x)={}&\log\!\sum_{t\in Y_{\mathrm{word}}}p_{\mathcal M}(t\mid u_j(c,x))\\
&-\log\!\sum_{t\in N_{\mathrm{word}}}p_{\mathcal M}(t\mid u_j(c,x)).
\end{split}
\label{eq:wordingscore}
\end{equation}
The affirmative set contains the distinct token identities of
\emph{yes}, \emph{Yes}, their leading-space forms, and \emph{YES},
retaining a form only when the complete form is a single token; the
negative set uses the analogous forms of \emph{no}. The uppercase
forms have no leading space in this scorer. All five wordings use
these same sets, the answer-position logits, and thinking-disabled
prompts. The main report instead uses the first token of each of six
case-and-space forms, as defined in App.~\ref{app:prompts}. The two
aggregations account for the Mistral measurement difference in
App.~\ref{app:conventions}.

The exact question for each wording follows the passage. In the quoted
instructions, \emph{candidate} is replaced by the candidate concept.

\begin{sdabox}{Passage only}
Based only on the passage above, is the concept \emph{candidate} one
of the most important things to remember in order to answer possible
future questions? Answer with a single word: yes or no.
\end{sdabox}

\begin{sdabox}{Unrestricted}
Is the concept \emph{candidate} one of the most important things to
remember in order to answer possible future questions about the passage
above? Answer with a single word: yes or no.
\end{sdabox}

\begin{sdabox}{Implied-content wording}
Considering both what the passage above states and what it implies
without stating, is the concept \emph{candidate} one of the most
important things to remember in order to answer possible future
questions? Answer with a single word: yes or no.
\end{sdabox}

\begin{sdabox}{Provenance-cued wording}
Would storing the concept \emph{candidate} help answer a likely future
question about the passage above, even if the passage never mentions
it? Answer with a single word: yes or no.
\end{sdabox}

\begin{sdabox}{Two-concept budget}
A memory system can keep only two concepts from the passage above.
Should the concept \emph{candidate} be one of them, so that a later
question can be answered? Answer with a single word: yes or no.
\end{sdabox}

\subsection{Matched naming and position measurements}
\label{app:naming}

Naming the useful concept raises both internal readouts more than
either importance report on Qwen3-8B, with episodes, candidates,
labels and questions fixed. Workspace, output-likelihood, report and rating AUCs
are initially 0.543, 0.613, 0.322 and 0.568. Appending a sentence naming
the useful concept raises them to 0.952, 0.859, 0.515 and 0.581;
placing the same sentence before the passage yields 0.916, 0.850,
0.582 and 0.655. Naming a distractor instead yields 0.393, 0.520,
0.249 and 0.519. Across six models, the stated-bridge control raises
mean workspace AUC from 0.563 to 0.956 and mean report AUC from 0.373
to 0.504 (Table~\ref{tab:statedcontrol}). The six individual report
changes include Mistral-7B's change from 0.411 to 0.395.

\begin{table}[tbp]
\centering
\caption{Matched stated-bridge control on Fresh-2. Appending a sentence naming the useful concept keeps candidates, labels and question fixed. Entries are pooled AUC. $W$: workspace readout; $\mathrm{LL}$: output likelihood; Report: yes/no importance report; Rating: 1--10 importance rating.}
\label{tab:statedcontrol}
\small
\setlength{\tabcolsep}{4pt}
\begin{tabular}{lcccccccc}
\toprule
& \multicolumn{2}{c}{$W$} & \multicolumn{2}{c}{$\mathrm{LL}$} & \multicolumn{2}{c}{Report} & \multicolumn{2}{c}{Rating} \\
\cmidrule(lr){2-3}\cmidrule(lr){4-5}\cmidrule(lr){6-7}\cmidrule(lr){8-9}
Model & Inferred & Stated & Inferred & Stated & Inferred & Stated & Inferred & Stated \\
\midrule
Qwen3-8B & 0.543 & 0.952 & 0.613 & 0.859 & 0.322 & 0.515 & 0.568 & 0.581 \\
Qwen3-14B & 0.571 & 0.954 & 0.636 & 0.864 & 0.355 & 0.607 & 0.565 & 0.591 \\
Qwen3-32B & 0.603 & 0.979 & 0.654 & 0.878 & 0.398 & 0.638 & 0.530 & 0.510 \\
Qwen2.5-7B & 0.520 & 0.879 & 0.600 & 0.890 & 0.427 & 0.456 & 0.481 & 0.481 \\
OLMo-2-7B & 0.561 & 0.986 & 0.633 & 0.911 & 0.327 & 0.410 & 0.533 & 0.585 \\
Mistral-7B & 0.581 & 0.990 & 0.630 & 0.854 & 0.411 & 0.395 & 0.608 & 0.601 \\
\midrule
Mean & 0.563 & 0.956 & 0.628 & 0.876 & 0.373 & 0.504 & 0.548 & 0.558 \\
\bottomrule
\end{tabular}
\end{table}


At $k=2$, Qwen3-8B SDA recall on the stated-bridge control is 55.4\%,
55.6\% and 52.9\%, against 44.9\% for the report and 39.9\% for the
reversed report. Naming the bridge changes what an omission or reversal
rule selects, while the trained module continues to retain useful
information.

\section{Extended Ranking Measurements}
\label{app:measurements}

\subsection{Internal readouts and shuffled context}
\label{app:context}

The report ranks useful content above the workspace readout on Explicit;
the ordering reverses on Decoupled across all eight models
(Table~\ref{tab:appA1}). Workspace rankings depend on the paired passage: shuffled context lowers
workspace AUC across all six models while preserving candidates and
labels (Table~\ref{tab:readoutshuffle}). All report measurements use
Eq.~\eqref{eq:score} and the admission prompt in App.~\ref{app:prompts}.

\begin{table}[t]
\centering
\caption{Shuffled context on Fresh-2. Each passage is replaced by another episode's passage, with candidates and labels fixed. Entries are pooled AUC with the paired passage and with shuffled context.}
\label{tab:readoutshuffle}
\small
\setlength{\tabcolsep}{5pt}
\begin{tabular}{lcccccc}
\toprule
 & \multicolumn{2}{c}{Workspace readout} & \multicolumn{2}{c}{Output likelihood} & \multicolumn{2}{c}{Yes/no report} \\
\cmidrule(lr){2-3} \cmidrule(lr){4-5} \cmidrule(lr){6-7}
Model & Paired & Shuffled & Paired & Shuffled & Paired & Shuffled \\
\midrule
Qwen3-8B & 0.543 & 0.433 & 0.613 & 0.604 & 0.322 & 0.485 \\
Qwen3-14B & 0.571 & 0.473 & 0.636 & 0.625 & 0.355 & 0.404 \\
Qwen3-32B & 0.602 & 0.479 & 0.654 & 0.623 & 0.398 & 0.416 \\
Qwen2.5-7B-Instruct & 0.520 & 0.509 & 0.600 & 0.610 & 0.427 & 0.441 \\
OLMo-2-7B-Instruct & 0.561 & 0.477 & 0.633 & 0.617 & 0.327 & 0.412 \\
Mistral-7B-Instruct & 0.581 & 0.423 & 0.630 & 0.620 & 0.411 & 0.455 \\
\bottomrule
\end{tabular}
\end{table}


\paragraph{Internal signals and the trained module use the passage.}
With shuffled context, workspace AUC falls across all six models,
including 0.543 to 0.433 on Qwen3-8B and 0.581 to 0.423 on Mistral-7B
(Table~\ref{tab:readoutshuffle}). Qwen3-8B likelihood AUC changes from
0.613 to 0.604 and report AUC from 0.322 to 0.485. Likelihood retains
a context-independent candidate prior, while workspace ranking and the
report's preference for stated candidates depend on the paired passage. Replacing each
passage with another episode's also reduces SDA admission AUC from
0.682 to 0.526, 0.555 and 0.534 across seeds (within-episode AUC
0.540, 0.549 and 0.542), demonstrating a context-specific ranking.


\subsection{Breadth across checkpoints, families and stages}
\label{app:breadth}

Across the 23-checkpoint, four-benchmark grid in Table~\ref{tab:appA1},
workspace AUC exceeds report AUC on Decoupled and Compositional at
every checkpoint with a defined report, including base and post-trained
stages. An entry of ``---'' means that the tokenizer ships no chat
template, so the yes/no importance report is undefined there.

\input{tables/app_full_predictive_grid}

Workspace AUC exceeds report AUC at every measured Qwen3 size on
Decoupled, with non-monotonic paired changes between adjacent sizes
(Table~\ref{tab:scale}). From 8B to 32B, the paired change is $-0.0536$
(95\% CI $[-0.1352,+0.0279]$).

\begin{table}[H]
\centering
\caption{Corrected Qwen3 Decoupled scale sweep. The endpoint contrast
$\Delta(W{-}V)$ from 8B to 32B is $+0.1027$ with paired 95\% CI
$[+0.0093,+0.1944]$; adjacent changes are non-monotonic.}
\label{tab:scale}
\small
\setlength{\tabcolsep}{7pt}
\begin{tabular}{lrrr}
\toprule
Model & $V$ & $W$ & $W{-}V$ \\
\midrule
\updatedrow Qwen3-1.7B & 0.3174 & 0.5607 & +0.2433$^{*}$ \\
\updatedrow Qwen3-4B & 0.4862 & 0.5291 & +0.0429 \\
\updatedrow Qwen3-8B & 0.4020 & 0.6545 & +0.2525$^{*}$ \\
\updatedrow Qwen3-14B & 0.3480 & 0.6677 & +0.3196$^{*}$ \\
\updatedrow Qwen3-32B & 0.3366 & 0.6919 & +0.3553$^{*}$ \\
\bottomrule
\end{tabular}
\vspace{2pt}
\parbox{0.94\linewidth}{\scriptsize $^{*}$Paired 95\% CI on $W{-}V$ excludes zero.}
\end{table}


Workspace AUC exceeds report AUC on Decoupled and Compositional at
every released OLMo-2 post-training stage measured at 1B and 7B; the
separation persists through SFT, DPO where available, and RLVR
(Table~\ref{tab:appA1}).


\subsection{Alternate-author and entity-bridge benchmarks}
\label{app:replications}

The report's preference for stated candidates persists with a different
author model. Across six models on the alternate-author benchmark,
report AUC ranges from 0.220 to 0.384, alongside provenance AUC
0.650--0.803 (Table~\ref{tab:replications}). These measurements use
the same report definition as the primary batteries.

\begin{table*}[t]
\centering
\caption{Replications on the alternate-author benchmark and entity-bridge
benchmark. Pooled AUCs: $W$, workspace readout; $V$, thinking-disabled
yes/no importance report. $W-V$: paired 95\% whole-episode bootstrap
interval (4,000 draws); $^{*}$ marks intervals excluding zero.}
\label{tab:replications}
\small
\begin{tabular}{llccc}
\toprule
Benchmark & Model & $W$ & $V$ & $W-V$ [95\% CI] \\
\midrule
Alternate-author & Qwen2.5-7B & 0.490 & 0.384 & +0.106 [-0.009, +0.215] \\
 & Qwen3-8B & 0.515 & 0.220 & \textcolor{gain}{+0.295} [+0.197, +0.386]$^{*}$ \\
 & Qwen3-14B & 0.521 & 0.267 & \textcolor{gain}{+0.254} [+0.166, +0.340]$^{*}$ \\
 & Qwen3-32B & 0.522 & 0.275 & \textcolor{gain}{+0.247} [+0.134, +0.356]$^{*}$ \\
 & OLMo-2-7B & 0.547 & 0.353 & \textcolor{gain}{+0.194} [+0.093, +0.297]$^{*}$ \\
 & Mistral-7B & 0.507 & 0.294 & \textcolor{gain}{+0.213} [+0.128, +0.300]$^{*}$ \\
\addlinespace
Entity-bridge & Qwen2.5-0.5B & 0.569 & 0.270 & \textcolor{gain}{+0.299} [+0.211, +0.396]$^{*}$ \\
 & Qwen2.5-7B & 0.664 & 0.881 & \textcolor{loss}{-0.217} [-0.303, -0.133]$^{*}$ \\
 & Qwen3-8B & 0.642 & 0.719 & -0.076 [-0.189, +0.032] \\
\bottomrule
\end{tabular}
\end{table*}


\paragraph{Report success accompanies a reversal of provenance preference.}
On the entity-bridge benchmark, Qwen2.5-7B and Qwen3-8B report AUCs
reach 0.881 and 0.719, above workspace AUCs of 0.664 and 0.642
(Table~\ref{tab:replications}). Both reports prefer unstated candidates:
their provenance AUCs are 0.133 and 0.289. Shuffled context reduces
report AUC to 0.282 and 0.340. Qwen2.5-0.5B retains the usual
preference for stated candidates, with provenance AUC 0.734 and utility
AUC 0.270. The report reads the paired passage in both regimes; the
direction of its provenance preference changes with candidate
construction.



\subsection{Admission after distillation across scales}
\label{app:sdascale}

SDA raises development admission AUC at every measured Qwen3 scale
and seed (Table~\ref{tab:sdascale}). All 18 paired 95\% intervals for
the gain lie above zero, while the largest absolute workspace AUC
change is 0.0046. The measured improvement is concentrated in the
admission interface.

\begin{table}[tbp]
\centering
\caption{Admission ranking on the Decoupled and Compositional development batteries across Qwen3 scales. These batteries informed configuration choice and were excluded from gradient training. $V$: pooled admission AUC; $\Delta V$: paired change with a 95\% whole-episode bootstrap interval; $|\Delta W|$: largest absolute workspace change across seeds.}
\label{tab:sdascale}
\small
\setlength{\tabcolsep}{4pt}
\begin{tabular}{llccccc}
\toprule
Model & Benchmark & $V$ before & Seed & $V$ after & $\Delta V$ [95\% CI] & $|\Delta W|$ \\
\midrule
\multirow{3}{*}{Qwen3-8B} & \multirow{3}{*}{Decoupled} & \multirow{3}{*}{0.387} & 0 & \textbf{0.756} & $+0.369$ [$+0.310$, $+0.431$] & \multirow{3}{*}{0.0033} \\
 & & & 1 & \textbf{0.739} & $+0.352$ [$+0.293$, $+0.417$] & \\
 & & & 2 & \textbf{0.758} & $+0.371$ [$+0.312$, $+0.435$] & \\
\addlinespace
\multirow{3}{*}{Qwen3-8B} & \multirow{3}{*}{Compositional} & \multirow{3}{*}{0.208} & 0 & \textbf{0.618} & $+0.410$ [$+0.325$, $+0.490$] & \multirow{3}{*}{0.0021} \\
 & & & 1 & \textbf{0.623} & $+0.415$ [$+0.329$, $+0.494$] & \\
 & & & 2 & \textbf{0.619} & $+0.412$ [$+0.326$, $+0.491$] & \\
\addlinespace
\multirow{3}{*}{Qwen3-14B} & \multirow{3}{*}{Decoupled} & \multirow{3}{*}{0.427} & 0 & \textbf{0.702} & $+0.275$ [$+0.216$, $+0.340$] & \multirow{3}{*}{0.0029} \\
 & & & 1 & \textbf{0.669} & $+0.242$ [$+0.192$, $+0.299$] & \\
 & & & 2 & \textbf{0.680} & $+0.253$ [$+0.200$, $+0.312$] & \\
\addlinespace
\multirow{3}{*}{Qwen3-14B} & \multirow{3}{*}{Compositional} & \multirow{3}{*}{0.216} & 0 & \textbf{0.529} & $+0.312$ [$+0.234$, $+0.395$] & \multirow{3}{*}{0.0029} \\
 & & & 1 & \textbf{0.458} & $+0.242$ [$+0.172$, $+0.317$] & \\
 & & & 2 & \textbf{0.481} & $+0.265$ [$+0.190$, $+0.346$] & \\
\addlinespace
\multirow{3}{*}{Qwen3-32B} & \multirow{3}{*}{Decoupled} & \multirow{3}{*}{0.478} & 0 & \textbf{0.773} & $+0.295$ [$+0.237$, $+0.356$] & \multirow{3}{*}{0.0037} \\
 & & & 1 & \textbf{0.762} & $+0.284$ [$+0.224$, $+0.346$] & \\
 & & & 2 & \textbf{0.721} & $+0.243$ [$+0.186$, $+0.299$] & \\
\addlinespace
\multirow{3}{*}{Qwen3-32B} & \multirow{3}{*}{Compositional} & \multirow{3}{*}{0.277} & 0 & \textbf{0.622} & $+0.344$ [$+0.273$, $+0.418$] & \multirow{3}{*}{0.0046} \\
 & & & 1 & \textbf{0.624} & $+0.347$ [$+0.267$, $+0.427$] & \\
 & & & 2 & \textbf{0.588} & $+0.310$ [$+0.233$, $+0.388$] & \\
\bottomrule
\end{tabular}

\end{table}


\section{Teacher and Optimization Measurements}
\label{app:ablations}

\subsection{Teacher composition and candidate readout}

\label{app:teacherreadout}

Table~\ref{tab:teachermodels} reports primary-pool recall for all three
teacher compositions; Table~\ref{tab:teachermodelpaired} gives their
complete paired intervals.

\begin{table}[htbp]
\centering
\caption{Teacher composition on the statedness-balanced pool for five models. Recall in percent, three training seeds per teacher; every other setting is held fixed. Each reduced column trains the same adapter on labels from SDA's teacher with one signal removed. Table~\ref{tab:teachermodelpaired} gives paired differences with 95\% intervals.}
\label{tab:teachermodels}
\small
\setlength{\tabcolsep}{4pt}
\begin{tabular}{llccc}
\toprule
Model & Budget & \textbf{SDA} & \shortstack{Without the\\workspace readout} & \shortstack{Without the\\output likelihood} \\
\midrule
Qwen3-8B & $k=1$ & \textbf{40.5} & 42.1 & 36.7 \\
 & $k=2$ & \textbf{53.7} & 54.8 & 49.8 \\
 & $k=3$ & \textbf{51.8} & 50.6 & 50.5 \\
\addlinespace
Qwen3-14B & $k=1$ & \textbf{35.5} & 36.2 & 34.3 \\
 & $k=2$ & \textbf{44.1} & 43.4 & 42.4 \\
 & $k=3$ & \textbf{45.3} & 45.8 & 45.2 \\
\addlinespace
Qwen3-32B & $k=1$ & \textbf{44.9} & 43.4 & 45.7 \\
 & $k=2$ & \textbf{53.3} & 52.2 & 54.7 \\
 & $k=3$ & \textbf{54.8} & 55.3 & 56.6 \\
\addlinespace
Qwen2.5-7B & $k=1$ & \textbf{17.4} & 16.2 & 17.6 \\
 & $k=2$ & \textbf{33.3} & 32.2 & 32.4 \\
 & $k=3$ & \textbf{39.1} & 37.6 & 37.3 \\
\addlinespace
OLMo-2-7B & $k=1$ & \textbf{49.5} & 46.2 & 45.2 \\
 & $k=2$ & \textbf{54.7} & 53.4 & 51.1 \\
 & $k=3$ & \textbf{55.5} & 55.5 & 55.2 \\
\bottomrule
\end{tabular}
\end{table}


\begin{table}[htbp]
\centering
\caption{SDA minus the student of the teacher with one signal removed on the statedness-balanced pool, for each of five models. Differences in recall are seed-averaged points with 95\% intervals from 4,000 paired whole-episode bootstrap draws; three seeds per teacher and model.}
\label{tab:teachermodelpaired}
\small
\setlength{\tabcolsep}{4pt}
\begin{tabular}{llccc}
\toprule
Model & Teacher & $k=1$ & $k=2$ & $k=3$ \\
\midrule
Qwen3-8B & Without workspace readout & $-1.6$ [$-4.0$, $+0.5$] & $-1.1$ [$-3.7$, $+1.6$] & $+1.2$ [$-1.2$, $+3.6$] \\
 & Without output likelihood & $+3.8$ [$+1.6$, $+6.2$] & $+4.0$ [$+1.9$, $+6.1$] & $+1.3$ [$-0.6$, $+3.3$] \\
Qwen3-14B & Without workspace readout & $-0.7$ [$-2.7$, $+1.2$] & $+0.6$ [$-0.6$, $+2.0$] & $-0.6$ [$-2.0$, $+0.9$] \\
 & Without output likelihood & $+1.1$ [$-0.2$, $+2.5$] & $+1.6$ [$+0.1$, $+3.4$] & $+0.1$ [$-1.4$, $+1.6$] \\
Qwen3-32B & Without workspace readout & $+1.5$ [$+0.1$, $+2.9$] & $+1.0$ [$-0.1$, $+2.2$] & $-0.5$ [$-1.8$, $+1.0$] \\
 & Without output likelihood & $-0.8$ [$-2.9$, $+1.1$] & $-1.5$ [$-3.1$, $+0.2$] & $-1.7$ [$-3.6$, $+0.0$] \\
Qwen2.5-7B & Without workspace readout & $+1.1$ [$-0.5$, $+2.6$] & $+1.1$ [$-0.6$, $+2.9$] & $+1.5$ [$+0.0$, $+2.9$] \\
 & Without output likelihood & $-0.3$ [$-1.6$, $+1.1$] & $+0.9$ [$-1.1$, $+2.9$] & $+1.8$ [$+0.3$, $+3.5$] \\
OLMo-2-7B & Without workspace readout & $+3.3$ [$+1.2$, $+5.5$] & $+1.3$ [$-0.5$, $+3.2$] & $+0.0$ [$-1.9$, $+1.8$] \\
 & Without output likelihood & $+4.3$ [$+2.0$, $+6.7$] & $+3.6$ [$+1.2$, $+6.1$] & $+0.3$ [$-1.9$, $+2.5$] \\
\bottomrule
\end{tabular}
\end{table}


Table~\ref{tab:crosspairmodels} gives the paired-passage results for each model.

\begin{table}[tbp]
\centering
\caption{Paired-passage results for each model: 181 pairs, 362 episodes per model. Entries are SDA minus the student of the teacher with one signal removed, averaged over three seeds, in points. Recall is blind semantic answer grading; retention indicates admission of the episode's useful concept. The 95\% intervals use 4,000 bootstrap draws resampling pairs within each model.}
\label{tab:crosspairmodels}
\footnotesize
\setlength{\tabcolsep}{2pt}
\begin{tabular}{llcccc}
\toprule
& & \multicolumn{2}{c}{Without the workspace readout} & \multicolumn{2}{c}{Without the output likelihood} \\
\cmidrule(lr){3-4}\cmidrule(lr){5-6}
Model & Budget & Retention & Recall & Retention & Recall \\
\midrule
Qwen3-8B & $k=1$ & $+5.3$ [$+2.5$, $+8.2$] & $-0.1$ [$-2.0$, $+1.8$] & $+2.9$ [$+0.0$, $+5.8$] & $+3.0$ [$+1.1$, $+4.8$] \\
 & $k=2$ & $+3.1$ [$+0.2$, $+6.2$] & $-0.2$ [$-2.6$, $+2.1$] & $+2.4$ [$-0.2$, $+5.2$] & $+3.2$ [$+1.0$, $+5.5$] \\
 & $k=3$ & $+3.5$ [$+1.1$, $+6.0$] & $+1.7$ [$-0.6$, $+3.8$] & $+2.2$ [$+0.1$, $+4.3$] & $+2.3$ [$+0.3$, $+4.3$] \\
\addlinespace
Qwen3-14B & $k=1$ & $+9.0$ [$+6.5$, $+11.7$] & $+1.4$ [$-0.6$, $+3.5$] & $-1.4$ [$-3.5$, $+0.7$] & $+0.3$ [$-1.2$, $+1.8$] \\
 & $k=2$ & $+6.0$ [$+3.8$, $+8.5$] & $+0.0$ [$-1.7$, $+1.8$] & $-1.3$ [$-3.3$, $+0.7$] & $+1.3$ [$-0.5$, $+3.1$] \\
 & $k=3$ & $+2.3$ [$+0.6$, $+4.0$] & $+0.6$ [$-0.9$, $+2.0$] & $+1.6$ [$-0.1$, $+3.3$] & $+1.8$ [$+0.0$, $+3.5$] \\
\addlinespace
Qwen3-32B & $k=1$ & $+3.8$ [$+2.1$, $+5.6$] & $+1.6$ [$+0.2$, $+3.0$] & $-10.6$ [$-13.3$, $-7.9$] & $-1.9$ [$-4.0$, $+0.1$] \\
 & $k=2$ & $+1.8$ [$-0.1$, $+3.7$] & $+0.7$ [$-0.6$, $+2.0$] & $-6.2$ [$-8.0$, $-4.4$] & $-0.3$ [$-2.2$, $+1.7$] \\
 & $k=3$ & $+0.7$ [$-0.6$, $+2.2$] & $+0.2$ [$-0.9$, $+1.2$] & $-4.5$ [$-6.3$, $-2.9$] & $-0.8$ [$-2.3$, $+0.6$] \\
\addlinespace
Qwen2.5-7B & $k=1$ & $+8.3$ [$+6.3$, $+10.3$] & $+3.0$ [$+1.7$, $+4.5$] & $-1.6$ [$-4.0$, $+0.8$] & $+0.8$ [$-0.7$, $+2.5$] \\
 & $k=2$ & $+6.2$ [$+4.5$, $+7.9$] & $+1.8$ [$+0.5$, $+3.2$] & $+0.6$ [$-1.2$, $+2.4$] & $+0.6$ [$-1.6$, $+2.6$] \\
 & $k=3$ & $+4.4$ [$+2.8$, $+6.2$] & $+1.5$ [$-0.1$, $+3.1$] & $+4.3$ [$+2.6$, $+6.2$] & $+0.4$ [$-1.5$, $+2.1$] \\
\addlinespace
OLMo-2-7B & $k=1$ & $+10.0$ [$+7.6$, $+12.5$] & $+3.4$ [$+1.5$, $+5.4$] & $+3.9$ [$+0.8$, $+6.8$] & $+4.4$ [$+2.4$, $+6.5$] \\
 & $k=2$ & $+6.0$ [$+4.1$, $+7.9$] & $+0.7$ [$-1.0$, $+2.4$] & $+5.0$ [$+2.5$, $+7.5$] & $+3.1$ [$+0.6$, $+5.6$] \\
 & $k=3$ & $+4.0$ [$+2.5$, $+5.5$] & $+1.7$ [$+0.5$, $+2.9$] & $+4.8$ [$+2.6$, $+7.1$] & $+3.4$ [$+1.3$, $+5.5$] \\
\bottomrule
\end{tabular}
\end{table}


We evaluate 362 Qwen3-8B episodes in 181 pairs, each sharing an
identical candidate set across two passages and their own questions.
Each episode includes the other member's useful concept as a candidate,
so only the passage distinguishes which candidate is useful.
Table~\ref{tab:crosspair} reports SDA recall of 38.1/49.1/52.3\% at
one, two and three stored concepts, with 48.2\% retention at one.
Removing the workspace readout from SDA's teacher costs $5.3$ retention
points ($[+2.4,+8.4]$) and leaves recall unchanged on this set: the
difference is $-0.1$ points ($[-2.0,+1.8]$). Removing the output
likelihood costs both, $2.9$ retention points ($[+0.1,+5.7]$) and $3.0$
recall points ($[+1.1,+4.9]$). These differences are seed-averaged,
with 95\% bootstrap intervals over episodes from 4,000 draws. Removing
the workspace readout also costs $3.1$ and $3.5$ retention points at
two and three stored concepts, with both intervals excluding zero. The
both-of-a-pair column reports 20.8\% for SDA, 17.9\% for its teacher
without the workspace readout and 20.4\% for its teacher without the
output likelihood. A passage-blind candidate ordering cannot raise this
rate above zero at one stored concept. Against the untrained baselines
at one stored concept, SDA improves recall by $10.5$ points over the
yes/no importance report ($[+6.7,+14.3]$) and by $2.5$ points over the
1--10 importance rating ($[-1.0,+6.0]$).

\begin{table*}[t]
\centering
\caption{Results on the paired-passage battery: 362 Qwen3-8B episodes in
181 pairs. Recall: blind semantic answer grading, in percent.
Retention: percentage admitting the episode's own useful concept;
the last column gives the percentage of pairs where both members
retain their own useful concept. Rows named by a removed signal
are students of SDA's teacher with that signal removed. Trained
policies average three seeds; untrained rows score candidates directly.}
\label{tab:crosspair}
\small
\setlength{\tabcolsep}{5pt}
\begin{tabular}{lccccc}
\toprule
& \multicolumn{3}{c}{Recall} & \multicolumn{2}{c}{Retention} \\
\cmidrule(lr){2-4}\cmidrule(lr){5-6}
Admission policy & $k=1$ & $k=2$ & $k=3$ & $k=1$ & \shortstack{Both of\\a pair} \\
\midrule
\textbf{SDA} & \textbf{38.1} & \textbf{49.1} & \textbf{52.3} & \textbf{48.2} & \textbf{20.8} \\
\addlinespace
Without the workspace readout & 38.2 & 49.3 & 50.6 & 42.8 & 17.9 \\
Without the output likelihood & 35.2 & 45.9 & 50.0 & 45.3 & 20.4 \\
1--10 importance rating \citep{park2023generative} & 35.6 & 43.6 & 46.1 & 39.5 & 13.8 \\
Yes/no importance report & 27.6 & 38.7 & 43.4 & 18.8 & 3.9 \\
Output likelihood used directly & 26.8 & 39.5 & 40.6 & 21.6 & 3.3 \\
Workspace readout used directly & 22.6 & 34.0 & 37.3 & 17.1 & 5.0 \\
\bottomrule
\end{tabular}
\end{table*}


Over the 220 teacher-labeled episodes, each containing a useful
candidate, SDA's teacher's two positive labels differ from those of the
teacher without the workspace readout on 43.2\% of episodes and from
those of the teacher without the output likelihood on 45.0\%. On
disagreement episodes, SDA's teacher's labels contain the useful concept
in 63.2\% against 58.9\% for the teacher without the workspace readout
(exact McNemar $p=0.67$), and in 63.6\% against 48.5\% for the teacher
without the output likelihood ($p=0.032$). Across all 220 episodes, these
rates are 60.9\%, 59.1\% and 54.1\%, respectively.

The teacher uses each candidate's first token as specified in
Eq.~\eqref{eq:wrr}. The full-candidate readout protocol appends the candidate and reads
each token at the position preceding it. Each token receives its peak
reciprocal rank over layers; the candidate score aggregates these peaks
by their mean or minimum. Each aggregation and the first-token readout
take their own maximum over surface forms.

On the standard Fresh-2 pool with Qwen3-8B, 56\% of candidates span
multiple tokens. Averaging over all candidate tokens improves ranking
over the matched first-token readout (Table~\ref{tab:candidatereadout}).

\begin{table*}[t]
\centering\small
\caption{Candidate readout comparison on the standard pool, Qwen3-8B,
363 episodes and 1,607 candidates. Entries are pooled and
within-episode AUCs.}
\label{tab:candidatereadout}
\begin{tabular}{lcc}
\toprule
Workspace readout & Pooled AUC & Within-episode AUC \\
\midrule
First token & 0.530 & 0.536 \\
Mean across all tokens & 0.574 & 0.571 \\
Minimum across all tokens & 0.531 & 0.538 \\
\bottomrule
\end{tabular}

\end{table*}

\paragraph{Teacher composition on the standard pool.}
On the standard Fresh-2 pool (Table~\ref{tab:sdaablation}), the teacher without the workspace readout reaches 50.3\%, a paired
$+0.3$ percentage points relative to SDA (95\% CI $[-2.2,+2.8]$);
the teacher without the output likelihood reaches 47.9\%, a paired $-2.1$ points
($[-4.1,-0.3]$).

\begin{table}[tbp]
\centering
\caption{SDA ablations on the standard Fresh-2 pool (Qwen3-8B, $k=2$, blind semantic recall in percent). Trained variants average three seeds. $\Delta$: recall difference from SDA with a paired, seed-averaged 4,000-draw episode bootstrap 95\% interval. The untrained report and the workspace readout and output likelihood used directly are in Table~\ref{tab:baselines}.}
\label{tab:sdaablation}
\small
\setlength{\tabcolsep}{5pt}
\begin{tabular}{lcc}
\toprule
Variant & Recall & $\Delta$ vs.\ SDA \\
\midrule
\textbf{SDA} & \textbf{50.0} & \emph{reference} \\
\addlinespace
\multicolumn{3}{l}{\emph{Teacher with one signal removed}} \\
Without the output likelihood & 47.9 & $-2.1$ [$-4.1$, $-0.3$] \\
Without the workspace readout & 50.3 & $+0.3$ [$-2.2$, $+2.8$] \\
\addlinespace
\multicolumn{3}{l}{\emph{No distillation}} \\
Combined teacher used directly & 45.7 & $-4.3$ [$-8.9$, $+0.1$] \\
\addlinespace
\multicolumn{3}{l}{\emph{Recall reward}} \\
Recall reward instead of distillation & 46.4 & $-3.7$ [$-6.2$, $-1.3$] \\
SDA followed by recall-reward refinement & 52.9 & $+2.8$ [$+0.6$, $+5.1$] \\
\bottomrule
\end{tabular}
\end{table}


\paragraph{Distillation changes the usable ranking.}
Using the combined teacher directly yields 45.7\% standard-pool Fresh-2 recall,
$-4.3$ points relative to SDA ($[-8.9,+0.1]$). Direct workspace scoring
yields 44.6\%, a difference of $-5.4$ points ($[-10.4,-0.5]$), and
direct likelihood yields 48.8\%, a difference of $-1.3$ points
($[-5.6,+2.8]$). The student learns a context-and-candidate judgment
from the teacher's ordering.

\subsection{The training objective determines admission ordering}
\label{app:objective}

The objective baseline aligns training with the admission score by
applying binary cross-entropy to the yes/no log-odds in
Eq.~\eqref{eq:score}. Teacher labels, training episodes, adapter shape
and the 414-step budget are fixed. Validation within-episode admission
AUC starts at 0.657 for the untrained model and reaches 0.824 at
the final step under SDA's vocabulary-normalized token objective.
At learning rate $10^{-5}$, the log-odds objective ends at
0.250, 0.333 and 0.301 across the three seeds. Validation selects the
untrained checkpoint in all three runs (Table~\ref{tab:objective}).

At a step size five times smaller ($2\times10^{-6}$), validation
selects a trained checkpoint after 100 steps, with within-episode AUC
0.690 and standard-pool recall of 40.2\% at $k=2$, against 50.0\%
for SDA's vocabulary-normalized token prediction. At $5\times10^{-6}$,
selection returns the untrained checkpoint. Each smaller-step-size run
uses one seed. SDA's token objective yields 9.8 points more recall
than the trained log-odds checkpoint and 9.6 points more than the
untrained checkpoint selected at the other learning rates.

\begin{table}[tbp]
\centering
\caption{Objective ablation on the standard pool, Qwen3-8B, $k=2$, blind semantic recall in percent. The main objective rows average three seeds at learning rate $10^{-5}$; each smaller-rate row uses one seed. Validation selects the checkpoint in every run.}
\label{tab:objective}
\small
\setlength{\tabcolsep}{6pt}
\begin{tabular}{lc}
\toprule
Training objective & Recall \\
\midrule
\textbf{SDA: token-level cross-entropy} & \textbf{50.0} \\
Yes/no log-odds cross-entropy & 40.5 \\
Yes/no log-odds, learning rate $2\times10^{-6}$ & 40.2 \\
Yes/no log-odds, learning rate $5\times10^{-6}$ & 40.5 \\
\bottomrule
\end{tabular}
\end{table}


\paragraph{Reader reward adds task-specific supervision.}
SDA exceeds recall-reward training by 3.7 points on Fresh-2
($[+1.3,+6.2]$); recall-reward refinement after SDA adds 2.8 points
($[+0.6,+5.1]$). These ablations use the same
training episodes and adapter shape with a matched step budget
(App.~\ref{app:ablationprotocol}).
